\subsection{HOMOGENEOUS SETS}

DEFINITION 6. Let G be a group. An operation of G on a set E is called transitive if there exists an element \(x \in E\) whose orbit is E. A G-set E is called homogeneous if the operation of G on E is transitive.

It is also said that G operates transitively on E; or that E is a homogeneous set under G. It amounts to the same to say that E is non-empty and that, for all elements x and y of E, there exists an element \(\alpha \in G\) such that \(\alpha.x = y\) .

Example. If E is a G-set, each orbit of E, with the induced operation, is a homogeneous set under G.

Let G be a group and H a subgroup of G. Consider the set G/H of left cosets modulo H. The group G operates on G/H on the left by \((g, xH) \mapsto gxH\) . Let N be the normalizer of H. The group N operates on G/H on the right by \((xH, n) \mapsto xHn = xnH\) . This operation induces on H the trivial operation and hence, on passing to the quotient, N/H operates on G/H on the right. Let \(\phi: (N/H)^{0} \to S_{G/H}\) be the homomorphism corresponding to this operation.

PROPOSITION 5. With the above notation, G/H is a homogeneous G-set. The mapping \(\phi\) induces an isomorphism of \((\mathrm{N} / \mathrm{H})^0\) onto the group of automorphisms of the G-set G/H.

The orbit in G/H of the element \(\dot{e}=H\) is G/H, whence the first assertion. We now prove the second. If \(n\in N\) defines by right translation the identity mapping on G/H, then \(\dot{e}.n=\dot{e}\) , that is H.n=H, whence \(n\in H\) . Therefore

N/H operates faithfully on G/H on the right and \(\phi\) is injective. The left operation of G and right operation of N/H on G/H commute and hence the operators of N/H define G-morphisms of G/H into itself, which are necessarily G-automorphisms since they are bijective. Therefore \(\phi\) takes its values in the group \(\Phi\) of G-automorphisms of G/H. We show that the image of \(\phi\) is \(\Phi\) . Let \(f \in \Phi\) . By transporting the structure, the stabilizer of \(f(\dot{e})\) in G is equal to the stabilizer of \(\dot{e}\) and hence to H. Let \(n \in G\) be such that \(f(\dot{e}) = n\dot{e}\) . The stabilizer of \(n\dot{e}\) in G is \(nHn^{-1}\) (no. 2, Proposition 2), whence \(nHn^{-1} = H\) and \(n \in N\) . For every element xH of G/H, \(f(xH) = f(x.\dot{e}) = x.f(\dot{e}) = xnH = xHn\) and f coincides with the mapping defined by n.

Remarks. (1) Let G be a group, H a subgroup of G and \(\phi: \mathbf{G} \to \mathfrak{S}_{\mathbf{G}/\mathbf{H}}\) the homomorphism corresponding to the operation of G on G/H. The kernel of \(\phi\) is the intersection of the conjugates of H (no. 2, Proposition 2). It is also the largest normal subgroup contained in H (no. 3). In particular, G operates faithfully on G/H if and only if the intersection of the conjugates of H reduces to \(\{e\}\) .

\begin{enumerate}
\def\labelenumi{(\arabic{enumi})}
\setcounter{enumi}{1}
\tightlist
\item
  Let G be a group and H and K subgroups such that H is a normal subgroup of K. Then K/H operates on the G-set G/H on the right and the canonical mapping of G/H onto G/K defines on passing to the quotient a G-set isomorphism \((G/H)/(K/H) \rightarrow G/K\) (cf.~no. 4).
\end{enumerate}

PROPOSITION 6. Let G be a group, E a homogeneous G-set, \(a \in E\) , H the stabilizer of \(a\) and K a subgroup of G contained in H. There exists one and only one G-morphism \(f\) of G/K into E such that \(f(e.K) = a\) . If K = H, \(f\) is an isomorphism.

If \(f\) is a solution, then \(f(x, \mathbf{K}) = x.a\) for all \(x\) in \(\mathbf{G}\) , whence the uniqueness; we show the existence. The orbital mapping defined by \(a\) is compatible with the equivalence relation \(y \in x\mathbf{K}\) on \(\mathbf{G}\) . For, if \(y = xk, k \in \mathbf{K}\) , then

\[
y. a = x k. a = x. a.
\]

A mapping \(f\) is thus derived of \(G / K\) into \(H\) which satisfies \(f(x, \mathbf{K}) = x.a\) for all \(x\) in \(G\) . This mapping is a G-morphism and \(f(\mathbf{K}) = a\) . This mapping is surjective for its image is a non-empty stable subset of \(E\) . Suppose now that \(K = H\) and let us show that \(f\) is injective. If \(f(x, \mathbf{H}) = f(y, \mathbf{H})\) , then \(x.a = y.a\) , whence \(x^{-1}y.a = a\) and \(x^{-1}y \in H\) , whence \(x.H = y.H\) .

THEOREM 1. Let G be a group.

\begin{enumerate}
\def\labelenumi{(\alph{enumi})}
\item
  Every homogeneous G-set is isomorphic to a homogeneous G-set of the form G/H where H is a subgroup of G.
\item
  Let H and H' be two subgroups of G. The G-sets G/H and G/H' are isomorphic if and only if H and H' are conjugate.
\end{enumerate}

As a homogeneous G-set is non-empty, assertion (a) follows from Proposition 6. We show (b). Let \(f: \mathrm{G} / \mathrm{H} \to \mathrm{G} / \mathrm{H}'\) be a G-set isomorphism. The subgroup H is the stabilizer of H and hence, by transport of structure, the stabilizer of an element of G/H'. The subgroups H and H' are therefore conjugate (no. 2, Proposition 2). If H' = αHα⁻¹, H' is the stabilizer of the element α. H of G/H (no. 2, Proposition 2) and hence G/H' is isomorphic to G/H (Proposition 6).

Examples. (1) Let E be a non-empty set. The group \(S_{E}\) operates transitively on E. If x and y are two elements of E, the mapping \(\tau: E \to E\) such that \(\tau(x) = y\) , \(\tau(y) = x\) and \(\tau(z) = z\) for \(z \neq x\) , y, is a permutation of E. Let \(a \in E\) . The stabilizer of a is identified with \(S_{F}\) , where \(F = E - \{a\}\) . The homogeneous G-set E is thus isomorphic to \(S_{E}/S_{F}\) .

\begin{enumerate}
\def\labelenumi{(\arabic{enumi})}
\setcounter{enumi}{1}
\tightlist
\item
  Let \(\mathbf{E}\) be a set of \(n\) elements and \((p_i)_{i\in \mathbf{I}}\) a finite family of integers \(>0\) such that \(\sum_{i}p_{i} = n\) . Let \(\mathbf{X}\) be the set of partitions \((\mathrm{F}_i)_{i\in \mathbf{I}}\) of \(\mathbf{E}\) such that \(\operatorname {Card}(\mathrm{F}_i) = p_i\) for all \(i\) . The group \(\mathfrak{S}_{\mathbf{E}}\) operates transitively on \(\mathbf{X}\) . The stabilizer H of an element \((\mathrm{F}_i)_{i\in \mathbf{I}}\) of \(\mathbf{X}\) is canonically isomorphic to \(\prod_{i\in \mathbf{I}}\mathfrak{S}_{\mathbf{F}_i}\) and is hence of order \(\prod_{i\in \mathbf{I}}p_i!\) . Applying Theorem 1 and §4, no. 4, Corollary to Proposition 4, a new proof is obtained of the fact that
\end{enumerate}

\[
\operatorname{Card} (\mathbf {X}) = \frac {n !}{\prod_ {i \in I} p _ {i} !}.
\]

In particular, take \(\mathbf{I} = \{1,2,\dots,r\}\) , \(\mathbf{E} = \{1,2,\dots,n\}\) ,

\[
\mathrm{F} _ {i} = \left\{p _ {1} + \dots + p _ {i - 1} + 1, \dots , p _ {1} + \dots + p _ {i} \right\}
\]

for \(1 \leqslant i \leqslant r\) . Let \(S\) be the set of \(\tau \in \mathfrak{S}_{E}\) such that \(\tau|_{F_i}\) is increasing for \(1 \leqslant i \leqslant r\) . If \((G_1, \ldots, G_r) \in X\) there exists one and only one \(\tau \in S\) which maps \((F_1, \ldots, F_r)\) to \((G_1, \ldots, G_r)\) . In other words, each left coset of \(\mathfrak{S}_{E}\) modulo \(H\) meets \(S\) in one and only one point.

\begin{enumerate}
\def\labelenumi{(\arabic{enumi})}
\setcounter{enumi}{2}
\tightlist
\item
  *Let \(n\) be an integer \(\geqslant 1\) . The orthogonal group \(\mathbf{O}(n, \mathbf{R})\) operates transitively on the unit sphere \(\mathbf{S}_{n-1}\) in \(\mathbf{R}^n\) . The stabilizer of the point \((0, \ldots, 0, 1)\) is identified with the orthogonal group \(\mathbf{O}(n-1, \mathbf{R})\) . The homogeneous \(\mathbf{O}(n, \mathbf{R})\) -set \(\mathbf{S}_{n-1}\) is thus isomorphic to \(\mathbf{O}(n, \mathbf{R}) / \mathbf{O}(n-1, \mathbf{R})\) .*
\end{enumerate}

